% !TeX spellcheck = en_GB
%!TEX root = ../main/main.tex
In this section we introduce an extension to the semantics of CEL, namely we introduce a new operator using \textit{Allen interval algebra} \textit{overlap}. 
We then extend the formal computational model for evaluating CEL formulas and prove its correctness. 

We present then the \emph{overlap operator} for CEL as with the following definition:
\[ 
\begin{aligned}
	\sem{\varphi_1 \overlap \varphi_2}(\Stream)
	&= \{C_1 \cup C_2 \mid C_1 \in \sem{\varphi_1}(\Stream) 
	\land C_2 \in \sem{\varphi_2}(\Stream) \\
	&\quad \land \cstart(C_1) \leq \cstart(C_2) \leq \cend(C_1) \leq \cend(C_2)\}
\end{aligned}
\]
 
\begin{example} Suppose we have sensors in a park monitoring humidity ($H$) and temperature ($T$). Consider a stream $\Stream = e_1 e_2 e_3 e_4$ where the event types are $\type(e_1) = H$, $\type(e_2) = T$, $\type(e_3) = H$, and $\type(e_4) = T$. 

Let $\varphi_1 = H ; H$ be a formula detecting a non contiguous sequence of humidity readings, and $\varphi_2 = T ; T$ be a formula detecting a non contiguous sequence of temperature readings.
\begin{itemize}
	\item The set $\sem{\varphi_1}(\Stream)$ contains a complex event $C_1$ representing the humidity sequence at positions $1$ and $3$, with $\cinterval(C_1) = [1..3]$. Thus, $\cstart(C_1) = 1$ and $\cend(C_1) = 3$.
	\item The set $\sem{\varphi_2}(\Stream)$ contains a complex event $C_2$ representing the temperature sequence at positions $2$ and $4$, with $\cinterval(C_2) = [2..4]$. Thus, $\cstart(C_2) = 2$ and $\cend(C_2) = 4$.
\end{itemize}

Since the overlap condition $\cstart(C_1) \leq \cstart(C_2) \leq \cend(C_1) \leq \cend(C_2)$ holds (i.e., $1 \leq 2 \leq 3 \leq 4$), the complex event $C = C_1 \cup C_2$ is in the result of $\sem{\varphi_1 \overlap \varphi_2}(\Stream)$. This resulting complex event captures the period where a humidity rise was followed by a temperature rise that continued after the humidity stabilized.

\cristian{Buen ejemplo!}
\cristian{Aqui habla de la subida de humedad y de temperatura, pero antes no se menciona nada sobre $\varphi_1$ y $\varphi_2$. Esta bien?}
\end{example}

We denote by $\celplus{\overlapsimb}$ the \emph{extension of CEL} with the overlap operation. We can prove the following result:

\begin{theorem} \label{theo:mytheorem}
	For every $\celplus{\overlapsimb}$ formula $\varphi$ there exists a CEA $\cA_\varphi$ such that $\sem{\varphi}(\Stream) = \sem{\cA_\varphi}(\Stream)$ for every stream $\Stream$.
\end{theorem}
\cristian{Cambie la presentación del teorema.}


\begin{proof}
	We already know that for every (standard) CEL operator we can construct an equivalent CEA. So, we concentrate next in the construction of a CEA for the overlap operator. 
	Let $\varphi_1$ and $\varphi_2$ be formulas in CEL extended with the overlap operator. We show by induction that there exists a CEA $\cAoverlap$ such that:
	\[
	\sem{\varphi_1 \overlap \varphi_2}(\Stream) = \sem{\cAoverlap}(\Stream)
	\]
	Assume then that there exist CEAs that satisfy the previous property for $\varphi_1$ and $\varphi_2$, called $\cA_{\varphi_1} = (Q_1, \SETP_1, \SETX_1, \Delta_1, q_0, F_1)$ and $\cA_{\varphi_2} = (Q_2, \SETP_2, \SETX_2, \Delta_2, p_0, F_2)$, respectively. 
	Then the construction for the overlap operator is as follows.
	Let $\cAoverlap$ be a CEA where $\cAoverlap = (\Qoverlap, \Poverlap, \Xoverlap, \Doverlap, q_0, F_2)$:
	\[
	\begin{aligned}[t]
		& \Qoverlap = Q_1 \uplus (Q_1 \times Q_2 \times \{0, 1\}) \uplus Q_2 \\
		& \Poverlap = \SETP_1 \cup \SETP_2 \cup \{ P_1 \land P_2 \mid P_1 \in \SETP_1 \land P_2 \in \SETP_2 \} \hspace*{155px}\\
		& \Xoverlap  = \SETX_1 \cup \SETX_2 \\
		& q_0 \text{ the initial state of }\cA_{\varphi_1}.\\
		& F_2 \text{ the final states of } \cA_{\varphi_2}.
	\end{aligned}
	\]
	The transition relation $\Doverlap$ is formally defined as:
	\[
	\begin{aligned}[t]
		\Doverlap & = \Delta_1 \uplus \Delta_2 \\
		& \uplus \{(q, \epsilon, (q, p_0, 0)) \mid q \in Q_1\} \\
		& \uplus \{((q, p, b), P_1 \land P_2, L_1 \cup L_2, (q', p', 1)) \mid (q, P_1, L_1, q') \in \Delta_1 \land (p, P_2, L_2, p') \in \Delta_2 \land b \in \{0, 1\}\} \\
		& \uplus \{((q, p, b), \epsilon, (q', p, b)) \mid (q, \epsilon, q') \in \Delta_1 \land p \in Q_2 \land b \in \{0, 1\}\}  \\
		& \uplus \{((q, p, b), \epsilon, (q, p', b)) \mid q \in Q_1 \land (p, \epsilon, p') \in \Delta_2 \land b \in \{0, 1\}\}  \\
		& \uplus \{((q, p, 1), \epsilon, p) \mid q \in F_1 \land p \in Q_2\}
	\end{aligned}
	\]\cristian{Revisado hasta aqui.}
Intuitively, given an stream S we capture the eventes given $\varphi_1$, and at some point (the overlap) we start capturing the events for $\varphi_2$ too.

Next, we prove that the construction is correct, namely, $\sem{\varphi_1 \overlap \varphi_2}(\Stream) = \sem{\cAoverlap}(\Stream)$ for every stream $\Stream$. The proof is by double containment. 

\paragraph{$\sem{\varphi_1 \overlap \varphi_2}(\Stream) \subseteq \sem{\cAoverlap}(\Stream)$}
Let $C_1 \cup C_2 \in \sem{\varphi_1 \overlap \varphi_2}(\Stream)$ where $C_1 \in \sem{\varphi_1}(\Stream)$ and $C_2 \in \sem{\varphi_2}(\Stream)$. We know there exist accepting runs $\rho_1$ of $\cA_{\varphi_1}$ and $\rho_2$ of $\cA_{\varphi_2}$ over $\Stream$ such that $C_{\rho_1} = C_1$ and $C_{\rho_2} = C_2$.
By the definition of a run, let:
\[ \rho_1 := (q_0, i_1) \xrightarrow{t_1} (q_1, i_2) \xrightarrow{t_2} \dots \xrightarrow{t_n} (q_{n}, i_{n+1}) \]
\[ \rho_2 := (p_0, j_1) \xrightarrow{t'_1} (p_1, j_2) \xrightarrow{t'_2} \dots \xrightarrow{t'_m} (p_{m}, j_{m+1}) \]
where $q_0$ and $p_0$ are the initial states, $q_{n} \in F_1$, and $p_{m} \in F_2$. Note that by definition, $\cstart(C_1) = i_1$, $\cend(C_1) = i_{n+1} - 1$, $\cstart(C_2) = j_1$, and $\cend(C_2) = j_{m+1} - 1$.

Because $C_1 \cup C_2$ is an overlapping event, we know $i_1 \leq j_1 \leq i_{n+1} - 1 \leq j_{m+1} - 1$.
To prove the containment, we construct the run $\rho_{\overlapsimb}$ of $\cAoverlap$ by combining $\rho_1$ and $\rho_2$. We define arbitrary indices $a \in [1..n]$ and $b \in [1..m]$ such that:
\begin{itemize}
	\item $a$ is the step in $\rho_1$ where the configuration first reaches stream position $j_1$ (i.e., $i_a = j_1$).
	\item $b$ is the step in $\rho_2$ where the configuration first reaches stream position $i_{n+1}$ (i.e., $j_b = i_{n+1}$).
\end{itemize}
We then move forward with the construction of $\rho_{\overlapsimb}$, using these points when needed.
First, the run $\rho_{\overlapsimb}$ mimics $\rho_1$ from step $1$ to $a-1$ using transitions in $\Delta_1$:
\[ 
\rho_{\overlapsimb} := (q_0, i_1) \xrightarrow{t_1} \dots \xrightarrow{t_{a-1}} (q_{a-1}, j_1)  \ldots 
\]
At this point we reach step $a$, here the stream is exactly at $j_1$, which is where $\rho_2$ must begin. We use the $\epsilon$-transition $(q_{a-1}, \epsilon, (q_{a-1}, p_0, 0)) \in \Doverlap$:
\[ \rho_{\overlapsimb} \ := \ \dots (q_{a-1}, j_1) \xrightarrow{\epsilon} ((q_{a-1}, p_0, 0), j_1) \dots \]

We now use both runs in the construction for the product state section of the transitions, carrying a bit $b$ that is initially $0$. We iterate over $x$ and $y$ where $x = a$ and $y = 1$. While $x \leq n$ and $y < b$, both runs are at the exact same stream position. At each step, we look at the transitions $t_x$ and $t'_y$ from $\rho_1$ and $\rho_2$ respectively:
\begin{itemize}
	\item If $t_x = \epsilon$, we use $\{((q_{x-1}, p_{y-1}, b), \epsilon, (q_{x}, p_{y-1}, b))\} \subseteq \Doverlap$. We advance $x \gets x+1$.
	\item If $t'_y = \epsilon$, we use $\{((q_{x-1}, p_{y-1}, b), \epsilon, (q_{x-1}, p_{y}, b))\} \subseteq \Doverlap$. We advance $y \gets y+1$.
	\item If $t_x = (P, L)$ and $t'_y = (P', L')$, both runs are reading the current stream event. We use $\{((q_{x-1}, p_{y-1}, b), P \land P', L \cup L', (q_{x}, p_{y}, 1))\} \subseteq \Doverlap$. We advance both $x \gets x+1$ and $y \gets y+1$, and the bit $b$ becomes $1$.
\end{itemize}
Because the overlap interval is valid ($j_1 \leq i_{n+1}-1$), there is at least one stream position where both automata consume an event. Thus, the bit $b$ is guaranteed to become $1$ during this phase.
By iterating these rules, the run reaches the point where $\rho_1$ is complete (we are in $q_n$) and $\rho_2$ is about to process step $b$ (we are in $p_{b-1}$):
\[ 
\rho_{\overlapsimb} \ := \ \dots ((q_{a-1}, p_0, 0), j_1) \rightarrow \overbrace{\dots}^{\text{iterative phase}} \rightarrow ((q_{n}, p_{b-1}, 1), i_{n+1})  \dots
\]
Since $q_{n} \in F_1$ and the bit is $1$, we can use the $\epsilon$-transition $((q_{n}, p_{b-1}, 1), \epsilon, p_{b-1}) \in \Doverlap$:
\[ 
\rho_{\overlapsimb} \ := \ \dots ((q_{n}, p_{b-1}, 1), i_{n+1}) \xrightarrow{\epsilon} (p_{b-1}, i_{n+1})  \dots 
\]
Finally, the run follows the remaining transitions of $\rho_2$ from step $b$ to $m$ using the transitions in $\Delta_2$:
\[ \rho_{\overlapsimb} \ := \ \dots (p_{b-1}, i_{n+1}) \xrightarrow{t'_b} \dots \xrightarrow{t'_m} (p_{m}, j_{m+1}) \]
Because $p_{m} \in F_2$, the run $\rho_{\overlapsimb}$ is accepting. Furthermore, by the construction of the iterative phase, every variable mapping $L$ captured by $\rho_1$ and $L'$ captured by $\rho_2$ are explicitly merged via $L \cup L'$. Thus, $C_{\rho_{\overlapsimb}} = C_1 \cup C_2$. Therefore, $\sem{\varphi_1 \overlap \varphi_2}(\Stream) \subseteq \sem{\cAoverlap}(\Stream)$.

\paragraph{$\sem{\cAoverlap}(\Stream) \subseteq \sem{\varphi_1 \overlap \varphi_2}(\Stream)$}
Let $C \in \sem{\cAoverlap}(\Stream)$. By definition, there exists an accepting run $\rho_{\overlapsimb}$ in $\cAoverlap$ over $\Stream$ such that $C_{\rho_{\overlapsimb}} = C$. 
Because $\rho_{\overlapsimb}$ starts at $q_0$ and ends at an accepting state $p_m \in F_2$, it must, by the construction of $\Doverlap$, contain exactly one entry $\epsilon$-transition $(q, \epsilon, (q, p_0, 0))$ and exactly one exit $\epsilon$-transition $((q', p, 1), \epsilon, p')$. Therefore, the run is strictly partitioned into three sections: one in $Q_1$, an overlap phase in $Q_1 \times Q_2 \times \{0,1\}$, and one in $Q_2$. Lets call them \emph{prefix}, \emph{overlap}, and \emph{suffix}, respectively.

We then can separate $\rho_{\overlapsimb}$ into two separate runs, $\rho_1$ and $\rho_2$, by isolating the components from the first and last partition.
\begin{itemize}
	\item \textbf{Construction of $\rho_1$:} For every state in the prefix and overlap phases of $\rho_{\overlapsimb}$, we take the first coordinate ($Q_1$). Because every transition in $\Doverlap$ that affects the first coordinate is strictly derived from $\Delta_1$, $\rho_1$ is a perfectly valid run in $\cA_{\varphi_1}$. Furthermore, because the overlap phase ends by taking $((q', p, 1), \epsilon, p')$, where $q' \in F_1$, $\rho_1$ is an accepting run. Let $C_1$ be the complex event generated by $\rho_1$.
	
	\item \textbf{Construction of $\rho_2$:} For every state in the overlap and suffix phases of $\rho_{\overlapsimb}$, we take the second coordinate ($Q_2$). Similarly, because the second coordinate is only modified by rules derived from $\Delta_2$, $\rho_2$ is a valid run in $\cA_{\varphi_2}$. Because $\rho_{\overlapsimb}$ ultimately ends in $F_2$, $\rho_2$ is an accepting run. Let $C_2$ be the complex event generated by $\rho_2$.
\end{itemize}
By induction, since $\rho_1$ and $\rho_2$ are accepting runs of $\cA_{\varphi_1}$ and $\cA_{\varphi_2}$, respectively, we have $C_1 \in \sem{\varphi_1}(\Stream)$ and $C_2 \in \sem{\varphi_2}(\Stream)$. 

Crucially, the exit transition out of the overlap phase is strictly conditioned on the state being $((q', p, 1), \epsilon, p')$. By the construction of $\Doverlap$, the only way the bit transitions from $0$ to $1$ is if $\rho_{\overlapsimb}$ executes at least one synchronized reading transition involving $P_1 \land P_2$ during the overlap phase. This strictly guarantees that there is at least one stream position $k$ where both $\rho_1$ and $\rho_2$ consume an event simultaneously. 
Consequently, $k \in \cinterval(C_1)$ and $k \in \cinterval(C_2)$. This implies that $\cstart(C_1) \leq k \leq \cend(C_1)$ and $\cstart(C_2) \leq k \leq \cend(C_2)$, from which it naturally follows that $\cstart(C_1) \leq \cstart(C_2) \leq \cend(C_1) \leq \cend(C_2)$. This directly resolves the interval constraints without ambiguity.

Because the overlapped transitions in $\Doverlap$ explicitly construct the variable mapping via $L_1 \cup L_2$, it follows that $C = C_1 \cup C_2$. Therefore, $C \in \sem{\varphi_1 \overlap \varphi_2}(\Stream)$.
\end{proof}

